% Beschreibung wie man vorgeht

\section{Einleitung}
Hier wird kurz zum Thema eingeleitet.
Sobald klar ist um was es geht, wird zu jedem Teil des Papers einmal gesagt was darin erarbeitet wird.
Im Abschnitt \ref{grundlagen} wird ein Überblick über den State-of-the-Art zum Thema xyz gegeben.
Daraufhin wird in Abschnitt ... usw.

\section{Grundlagen} \label{grundlagen}
Hier fangen die Grundlagen an.
Bitte mit einem einleitenden Satz anfangen und dann Related Work bzw. State-of-the-Art beschreiben.
Eine gute Vorgehensweise ist, einen Abschnitt pro Quelle/Paper zu nehmen.
Ein Abschnitt wird durch eine leere Zeile in Latex erzeugt.

Werkmeister \cite{dbt_mods_00050711} spricht in seiner Arbeit über Design Patterns und wendet dabei Entwicklungstechniken, die von Cooper et al. \cite{interactionDesign} beschrieben wurden.
Dies ist für das Thema xyz besonders relevant weil die Problemstellungen 1, 2 und 3 dabei herausgearbeitet wurden...\\

Mit dem Befehl \texttt{\textbackslash \textbackslash} am Ende einer Zeile und einer darauf folgenden leeren Zeile kriegt ihr einen Absatz mit einer freien Zeile dazwischen.
Setzt das sinnvoll, bedacht und konsequent ein.


\section{Hauptteil} \label{hauptteil}
Nun folgt der Hauptteil den ihr NICHT 'Hauptteil' nennt.
Hier sollte passend zu eurer Fragestellung / eurem Fokus die Einteilung des Texts geschehen.

\subsection{Subsection vom Hauptteil} \label{subsectionhauptteil}

Hallo ich bin eine Subsection von der Section  \ref{hauptteil}.

\subsubsection{Subsubsection... Level drei} 

Ich bin eine Subsubsection.
Ich bin ein Unterpunkt zu \ref{subsectionhauptteil}
Im Folgendem könnt ihr die Tabelle \ref{tab1} sehen.
Die Tabelle beinhaltet verschiedene Style-Optionen.
Da diese Tabelle zu breit für eine Spalte ist wurde statt \texttt{\textbackslash table} einfach ein * angehängt.
Also steht im Befehl \texttt{\textbackslash table*}.
\begin{table*}
\centering
\caption{Tabelle für Styles}\label{tab1}
\begin{tabular}{|l|l|l|}
\hline
Heading level &  Example & Font size and style\\
\hline
Title (centered) &  {\Large\bfseries Lecture Notes} & 14 point, bold\\
1st-level heading &  {\large\bfseries 1 Introduction} & 12 point, bold\\
2nd-level heading & {\bfseries 2.1 Printing Area} & 10 point, bold\\
3rd-level heading & {\bfseries Run-in Heading in Bold.} Text follows & 10 point, bold\\
4th-level heading & {\itshape Lowest Level Heading.} Text follows & 10 point, italic\\
\hline
\end{tabular}
\end{table*}

So könnte man eine mathematische Formel einbinden.
Das ist aber fürs Fachseminar wohl weniger interessant.

\begin{equation}
x + y = z
\end{equation}

Und Bilder werden am Besten als PDF oder in Vektorformaten eingebunden.
Ein Beispiel dafür wäre \ref{fig1}.

\begin{figure}[htbp]
\includegraphics[width=\columnwidth]{Resources/actorsdependencies2.pdf}
\caption{Zusammenspiel von Akteuren einer urbanen Datenplattform} \label{fig1}
\end{figure}

\subsection{Subsection vom Hauptteil die Zweite}
In der Bib-Datei habe ich euch noch ein paar nette Quellen hinterlassen.
Da könnt ihr auch sehen wie z.B. Websiten als Quellen angegeben werden.
Hierfür habe ich eine Statistik \cite{statistaVRinteresse} über das Interesse an VR in Deutschalnd raus gesucht.
DIN-Spezifikationen kann man auch als Quelle angeben.
Ich habe dazu kein offizielles Format gefunden, aber das Beispiel für \cite{DINSPEC91357} sollte formell keine Fehler enthalten.
Falls ein Artikel mehrere Autoren hat, wird der erste einfach genannt.
Wie bei dem Artikel von Huang et al. \cite{huang2015personal}.